\documentclass[12pt]{article}

% ============================================================
% PACKAGES
% ============================================================

\usepackage{amsmath}      % Mathematical environments: align, equation, etc.
\usepackage{amsthm}       % Theorem, definition, proof environments
\usepackage{amsfonts}     % Mathematical fonts such as \mathbb
\usepackage{amssymb}      % Additional mathematical symbols
\usepackage{mathtools}    % Extensions to amsmath
\usepackage{geometry}     % Page margins
\usepackage{enumitem}     % Customize enumerate/itemize
\usepackage{mdframed}     % Framed boxes
\usepackage{tikz}         % Draw diagrams
\usepackage{hyperref}     % Clickable references and links

\geometry{
    letterpaper,
    margin=0.8in
}

\linespread{1.2}


% ============================================================
% NUMBER SYSTEMS
% ============================================================
% Instead of writing \mathbb{R} every time, use \RR.

\newcommand{\RR}{\mathbb{R}}   % Real numbers
\newcommand{\CC}{\mathbb{C}}   % Complex numbers
\newcommand{\NN}{\mathbb{N}}   % Natural numbers
\newcommand{\ZZ}{\mathbb{Z}}   % Integers
\newcommand{\QQ}{\mathbb{Q}}   % Rational numbers


% ============================================================
% COMMON MATHEMATICAL NOTATION
% ============================================================
% These commands are useful for frequently used notation.

\newcommand{\set}[1]{\left\{#1\right\}}
% Example:
%   \set{x \in \RR : x > 0}
% produces:
%   { x in R : x > 0 }

\newcommand{\abs}[1]{\left|#1\right|}
% Example:
%   \abs{x-y}
% produces:
%   |x-y|

\newcommand{\norm}[1]{\left\|#1\right\|}
% Example:
%   \norm{x}
% produces:
%   ||x||

\newcommand{\inner}[2]{\left\langle #1,#2\right\rangle}
% Example:
%   \inner{x}{y}
% produces:
%   <x,y>

\newcommand{\parens}[1]{\left(#1\right)}
% Example:
%   \parens{x+y}

\newcommand{\brac}[1]{\left[#1\right]}
% Example:
%   \brac{a,b}

\newcommand{\angles}[1]{\left\langle #1\right\rangle}
% Example:
%   \angles{x,y}


% ============================================================
% THEOREM / THEORY COMPONENTS
% ============================================================

% "definition" style is upright text.
% Use this for definitions, examples, and remarks.

\theoremstyle{definition}

% ------------------------------------------------------------
% DEFINITION
% ------------------------------------------------------------
% Use when introducing a formal mathematical definition.
%
% Usage:
%
% \begin{definition}[Metric]
% A metric on a set X is a function ...
% \end{definition}
%
% [Metric] is optional and gives the definition a name.

\newtheorem{definition}{Definition}[section]


% ------------------------------------------------------------
% EXAMPLE
% ------------------------------------------------------------
% Use to give a concrete example of a definition/theorem.
%
% Usage:
%
% \begin{example}
% Consider X = \RR with the usual distance.
% \end{example}

\newtheorem{example}[definition]{Example}


% ------------------------------------------------------------
% REMARK
% ------------------------------------------------------------
% Use for observations, comments, or small additional facts.
% It is NOT intended for major mathematical results.
%
% Usage:
%
% \begin{remark}
% This definition does not require X to be a subset of \RR^n.
% \end{remark}

\newtheorem{remark}[definition]{Remark}


% ============================================================
% THEOREM-LIKE COMPONENTS
% ============================================================

% "plain" style makes the mathematical statements italic.

\theoremstyle{plain}


% ------------------------------------------------------------
% THEOREM
% ------------------------------------------------------------
% Use for an important/general mathematical result.
%
% Usage:
%
% \begin{theorem}[Cauchy--Schwarz Inequality]
% For all x,y \in \RR^n,
% \[
% |\inner{x}{y}| \leq \norm{x}\norm{y}.
% \]
% \end{theorem}

\newtheorem{theorem}[definition]{Theorem}


% ------------------------------------------------------------
% PROPOSITION
% ------------------------------------------------------------
% Use for a useful mathematical result that is usually less
% central than a theorem.
%
% Usage:
%
% \begin{proposition}
% Every finite-dimensional subspace of \RR^n is closed.
% \end{proposition}

\newtheorem{proposition}[definition]{Proposition}


% ------------------------------------------------------------
% LEMMA
% ------------------------------------------------------------
% Use for an auxiliary result whose main purpose is to help
% prove another theorem/proposition.
%
% Usage:
%
% \begin{lemma}
% ...
% \end{lemma}

\newtheorem{lemma}[definition]{Lemma}




% ------------------------------------------------------------
% COROLLARY
% ------------------------------------------------------------
% Use for a result that follows relatively directly from a
% theorem/proposition that was just established.
%
% Usage:
%
% \begin{corollary}
% ...
% \end{corollary}

\newtheorem{corollary}[definition]{Corollary}


% ============================================================
% PROOF
% ============================================================
% The proof environment comes from amsthm.
%
% Usage:
%
% \begin{proof}
% Suppose that ...
%
% Therefore, ...
% \end{proof}
%
% LaTeX automatically adds the QED symbol at the end.


% ============================================================
% CUSTOM BOX: IDEA
% ============================================================
% Use this when you want to record intuition or the main idea
% behind a mathematical concept.
%
% Usage:
%
% \begin{idea}
% \textbf{Key idea.}
% A metric gives us a notion of distance...
% \end{idea}

\newmdenv[
    linewidth=1pt,
    skipabove=10pt,
    skipbelow=10pt
]{idea}


% ============================================================
% CUSTOM BOX: WARNING
% ============================================================
% Use this for common mistakes, distinctions, or things that
% are easy to confuse.
%
% Usage:
%
% \begin{warning}
% Do not confuse continuity with uniform continuity.
% \end{warning}

\newmdenv[
    linewidth=1pt,
    skipabove=10pt,
    skipbelow=10pt
]{warning}

% ============================================================
% CUSTOM BOX: NOTE
% ============================================================
% Use this for a short personal note, useful observation,
% connection between concepts, or explanation.
%
% Unlike "Remark", this is intended for informal notes.
%
% Usage:
%
% \begin{note}
% A sequence can be viewed as a function from $\NN$ to $\RR$.
% \end{note}

\newmdenv[
    linewidth=1pt,
    skipabove=10pt,
    skipbelow=10pt
]{note}


% ============================================================
% DOCUMENT
% ============================================================

\begin{document}


% ============================================================
% TITLE
% ============================================================

\begin{center}
    {\Large \textbf{Series}}\\[0.5em]
    {\large Lecture 18}
\end{center}

\tableofcontents

\newpage

\section{Series}
\begin{definition}[Series]
Given $\{a_n\}$, let $\sum_{n=p}^qa_n = a_p+a_{p+1}+...+a_{q}$. Let $s_n = \sum_{p=1}^na_p$. Then $\{s_n\}$ is a sequence, sometime written as $\sum_{n=1}^{\infty}a_n$, called an infinite series.
\end{definition}

\begin{note}
If the defined sequence above converges to some $s$, we write: $\sum_{n=1}^{\infty}a_n=s$.
\end{note}

\section{Convergence of a series}
\begin{note}
Note that, we are dealing with sequences in $\RR$, then, if the sequence is Cauchy, then it converges.
\end{note}

\begin{example}
The harmonic series $\frac{1}{1}+\frac{1}{2}+\frac{1}{3}+...$ doesn't converge because it is not Cauchy.
\end{example}

\begin{theorem}[Cauchy criterion for series in $\RR$]
$\sum a_n$ converges $\iff$ for all $\epsilon > 0$ there is a $N$ such that $m,n\ge N$ then $\abs{\sum_{k=n}^ma_k}<\epsilon$.
\end{theorem}

\begin{note}
Note that, the right hand side of the theorem is the condition for a sequence is Cauchy, with the sequence is constructed by using Definition 1.1.
\end{note}

\begin{corollary}
$\sum a_n$ converges if $\lim_{n\rightarrow\infty}a_n = 0$.
\end{corollary}

\begin{note}
Note that, the series' elements are partial sums, and the $a_n$ in the $\lim$ of Corollary 2.3 is a single number in that partial sum. Corollary 2.3 is born simply by using $m=n$ at the right hand side.

Counter example for the inverse direction of Corollary 2.3: harmonic series.
\end{note}

\begin{theorem}[Convergence of non-negative series]
If $a_n \ge 0$ then $\sum a_n$ converges $\iff$ partial sums are bounded.
\end{theorem}

\begin{note}
Non-negative series in $\RR$ are monotonically increasing sequences. Then for those series to be converged, it must be bounded.
\end{note}

\begin{theorem}[Comparison test for convergence of series]
Here are some ways to check convergence of a series?
\begin{itemize}
    \item If $\abs{a_n}<c_n$ for large enough $n$, then if $\sum c_n$ converges, then $\sum a_n$ converges.
    \item $a_n\ge d_n\ge 0$ for large enough $n$, if $\sum d_n$ diverges, then $\sum a_n$ diverges.
\end{itemize}
\end{theorem}

\begin{note}
At Theorem 2.5, if you can prove the first comparison, then the next comparison is pretty easy by using contrapositive ($A \Rightarrow B \iff \overline{B} \Rightarrow \overline{A}$).

We have counter example for the first comparison if we don't have the absolute value operator at $a_n$: $a_n = \frac{-1}{n}$ which is diverged, $c_n = \frac{1}{n^2}$ which is converged.
\end{note}

\section{Geometric Series}
\begin{note}
In this part, I won't mention any proof or intuition.
\end{note}

\begin{theorem}[Convergence of Geometric Series]
If $\abs{x}<1$ then $\sum_{n=0}^{\infty}x^n$ converges, else diverges.
\end{theorem}

\begin{theorem}[Euler's number]
$\sum_{n=0}^{\infty}\frac{1}{n!}=e$.
\end{theorem}

\end{document}